\section{Catálogo de preguntas de decisión}\label{sec:comparacion:catalogo}\label{sec:hallazgos:catalogo-pi}

Una plantilla de pregunta es una pregunta con huecos, como la API, el plan, el volumen o el precio, que se rellenan para cada caso.
El catálogo tiene 32 plantillas, y 25 de ellas (el 78~\%) se pueden responder con el instrumental del canal API: la Datasheet, PRIME4API y HARVEY4API.
De esas 25, 11 se responden leyendo la Datasheet, 10 con una operación de PRIME4API y 4 con HARVEY4API.
Otras 2 se responden leyendo una matriz de \textit{features} por plan, que se captura con las mismas fuentes que la Datasheet pero queda fuera del esquema.
Las 5 restantes no tienen respuesta, y cada una deja escrito qué falta para tenerla.

La proporción de preguntas respondibles es la cobertura, la misma métrica con la que Cavero et al., que definen el Pricing Intelligence como disciplina, miden el alcance de sus herramientas: un 54~\% sobre un banco de 56 plantillas \citep{cavero2026pricingintelligence}.

Las dos cifras no se miden igual.
Cavero et al.\ prueban herramientas que ya existían contra un banco que no se hizo para ellas, y aquí la clase de cada pregunta, el instrumento con el que se responde, se asigna conociendo el instrumental.
La categoría B, la de selección de suscripción, parte de las operaciones del motor, y las cuatro preguntas que el catálogo asigna a HARVEY4API cuentan por la especificación de la herramienta, que no se ha ejecutado.
Sin esas cuatro, la cobertura con un instrumento leído o ejecutado es de 21 sobre 32, el 66~\%.

Del trabajo de Cavero et al.\ se toma también la gramática de las plantillas, con siete \textit{placeholders} (los huecos): la API (\texttt{SAAS}), el plan con sus complementos (\texttt{CONFIGURATION}), una funcionalidad (\texttt{FEATURE}), una cantidad como un volumen de correos o una ventana de tiempo (\texttt{USAGE LIMIT}), la necesidad del usuario (\texttt{ACTION}), un precio (\texttt{PRICE}) y el año de una versión del \textit{pricing} (\texttt{YEAR}).
Ninguna plantilla nombra una API, porque el proveedor solo aparece al rellenar los huecos, en el ejemplo de cada pregunta, y 13 de las 32 hablan de correo, como los envíos al mes o los límites que crecen con la reputación del remitente.
Con esa gramática, cada pregunta del catálogo puede incorporarse al banco de plantillas del grupo ISA de la Universidad de Sevilla.

Cada pregunta tiene una plantilla, un ejemplo instanciado sobre una de las ocho APIs y una clase, que dice con qué se responde.
Instanciar es rellenar los huecos: con SendGrid y 100.000, la plantilla de PI-B01, «suscripción más barata de \texttt{SAAS} para \texttt{USAGE LIMIT} correos al mes», da «la suscripción más barata de SendGrid para 100.000 correos al mes», una pregunta del caso de los 100.000 correos.
Los ejemplos instanciados están en \texttt{03\_datos/preguntas/catalogo.yaml}, en el repositorio del trabajo.

La clase D indica que basta leer la Datasheet, incluido su \textit{provenance}, las notas en las que la Datasheet guarda lo que los campos de ACM2YAML v0.3 no admiten.
Es el caso de los volúmenes contratados del plan Basic de Postmark, que se vende en seis volúmenes, cada uno con su precio.

La clase P indica que PRIME4API resuelve la pregunta con la operación que responde una de las diez preguntas de análisis de Gavira Sánchez \citep{gavira2026tfm}, y la tabla anota cuál.
El catálogo emplea siete de esas operaciones: el tiempo mínimo, la capacidad en una ventana, la suscripción más barata, la suscripción con presupuesto, el importe del sobrecoste, la demanda sostenible y las cuotas que se agotan a la vez.
Su correspondencia con el cálculo a mano de la restricción activa (el límite que se alcanza primero y frena el envío), el tiempo mínimo y el factor $k$ (los destinatarios por petición) está en la Sección~\ref{sec:limites:prime4api}.

Las otras tres clases son H, L y X.
La clase H pide combinar varias Datasheets o varias operaciones, que es la tarea de HARVEY4API \citep{ruizlopez2026tfm}.
La clase L se responde leyendo la matriz de \textit{features}, sin instrumento.
La clase X marca una pregunta sin respuesta y dice qué falta: un campo que v0.3 no serializa, una operación que el motor no tiene o un comportamiento que ninguna dimensión de capacidad del ACM recoge.
Una dimensión de capacidad es un aspecto del consumo que el proveedor limita por separado, como el volumen al mes o las peticiones por segundo.
La clase depende del instrumento y no de lo que publique cada proveedor, así que el reparto es el mismo en las ocho APIs: 25 preguntas respondibles con la Datasheet, PRIME4API o HARVEY4API, 2 que se responden leyendo la matriz y 5 sin respuesta.

Las 32 plantillas se agrupan en siete categorías, de la A a la G, según la fuente de la que salen.
La A recorre los bloques de capacidad de la Datasheet (cuota, caudal y sobrecoste) y la B parte de las operaciones de Gavira Sánchez sobre un plan.
La C compara APIs entre sí y la D convierte en preguntas la RQ3, la relación entre límites de uso y \textit{rates}.
Un límite de uso es el máximo de un nivel de uso medible, como la cuota de correos al mes.
La E, la F y la G salen de los elementos no soportados que se registraron al construir las Datasheets, los hechos del dominio que ACM2YAML v0.3 no puede expresar.
Recogen las necesidades funcionales, la vigencia de los precios y el riesgo de operar.

El catálogo se cerró en 32 preguntas sin un criterio de saturación fijado de antemano, y al pasar de tres APIs a ocho no ganó ninguna.
Tampoco cubre los cinco retos abiertos que enumeran Cavero et al.\ \citep{cavero2026pricingintelligence}, entre ellos la predicción de la evolución de un \textit{pricing} y los análisis \textit{what-if}.

Las tablas~\ref{tab:catalogo-pi-1}, \ref{tab:catalogo-pi-2} y~\ref{tab:catalogo-pi-3} recogen las 32 preguntas con su plantilla abreviada, su clase, la operación o composición que la resuelve y lo que falta para responder las X o cómo se leen las L.

\begin{table}[htbp]
\centering
\scriptsize
\begin{tabular}{|l|>{\raggedright\arraybackslash}p{6.6cm}|>{\raggedright\arraybackslash}p{2.2cm}|>{\raggedright\arraybackslash}p{3.4cm}|}
\hline
\textbf{Id} & \textbf{Pregunta (plantilla abreviada)} & \textbf{Clase y operación} & \textbf{Qué falta (X) o cómo se lee (L)} \\
\hline
\multicolumn{4}{|l|}{\textbf{A. Lectura de la capacidad contractual de una API (8)}} \\
\hline
PI-A01 & Correos al mes que incluye \texttt{CONFIGURATION} de \texttt{SAAS}, y qué pasa al superarlos & D & \\
PI-A02 & \textit{Rate} máximo con el que \texttt{SAAS} acepta peticiones de \texttt{ACTION} en \texttt{CONFIGURATION} & D & \\
PI-A03 & Destinatarios por petición de \texttt{ACTION} en \texttt{SAAS}, y si cuentan uno a uno contra la cuota & D & \\
PI-A04 & Si hay límite diario además del mensual en \texttt{CONFIGURATION} de \texttt{SAAS} & D & \\
PI-A05 & Tamaño máximo de un mensaje (adjuntos incluidos) en \texttt{SAAS} & D & \\
PI-A06 & Tiempo mínimo que necesita \texttt{CONFIGURATION} de \texttt{SAAS} para enviar \texttt{USAGE LIMIT} correos & P, tiempo mínimo & \\
PI-A07 & Capacidad que acumula \texttt{CONFIGURATION} de \texttt{SAAS} en una ventana de \texttt{USAGE LIMIT} & P, capacidad en una ventana & \\
PI-A08 & Peticiones simultáneas que admite \texttt{SAAS} en \texttt{CONFIGURATION} & X & Concurrencia: v0.3 no la serializa, y de las ocho APIs solo Mailgun publica algo sobre ella para la API REST: que no limita las peticiones simultáneas \\
\hline
\multicolumn{4}{|l|}{\textbf{B. Selección de suscripción para una necesidad (5)}} \\
\hline
PI-B01 & Suscripción más barata de \texttt{SAAS} para \texttt{USAGE LIMIT} correos al mes & P, suscripción más barata & \\
PI-B02 & Con \texttt{PRICE} al mes, máximo de correos con \texttt{SAAS} & P, suscripción con presupuesto & \\
PI-B03 & Con \texttt{USAGE LIMIT} al mes, si sale más barato el sobrecoste o el siguiente plan de \texttt{SAAS} & P, importe del sobrecoste & \\
PI-B04 & Si se puede empezar gratis en \texttt{SAAS} y qué limita el plan gratuito para \texttt{ACTION} & D & \\
PI-B05 & Plan de \texttt{SAAS} que sostiene un pico de \texttt{USAGE LIMIT} correos en \texttt{USAGE LIMIT 2} & P, demanda sostenible & \\
\hline
\end{tabular}
\caption[Catálogo de preguntas de decisión, categorías A y B]{Catálogo de preguntas de decisión, categorías A y B. Clase: D
Datasheet, P PRIME4API, H HARVEY4API, L lectura de la matriz de \textit{features}, X no respondible.}
\label{tab:catalogo-pi-1}
\end{table}

\begin{table}[htbp]
\centering
\scriptsize
\begin{tabular}{|l|>{\raggedright\arraybackslash}p{6.6cm}|>{\raggedright\arraybackslash}p{2.2cm}|>{\raggedright\arraybackslash}p{3.4cm}|}
\hline
\textbf{Id} & \textbf{Pregunta (plantilla abreviada)} & \textbf{Clase y operación} & \textbf{Qué falta (X) o cómo se lee (L)} \\
\hline
\multicolumn{4}{|l|}{\textbf{C. Comparación entre APIs (5)}} \\
\hline
PI-C01 & Para \texttt{USAGE LIMIT} correos al mes, cuál sale más barata, \texttt{SAAS 1} o \texttt{SAAS 2} & H, suscripción más barata por API, más comparación & \\
PI-C02 & Cuál de \texttt{SAAS 1}, \texttt{SAAS 2} y \texttt{SAAS 3} da más capacidad por dólar en el primer plan de pago & H, lectura por API, más comparación & \\
PI-C03 & Qué API de correo publica su \textit{rate} y cuál no & D & \\
PI-C04 & Cuál envía antes \texttt{USAGE LIMIT} correos, \texttt{SAAS 1} o \texttt{SAAS 2} & H, tiempo mínimo por API, más comparación & \\
PI-C05 & En qué se parecen y en qué se diferencian los pricings de \texttt{SAAS 1} y \texttt{SAAS 2} para \texttt{ACTION} & X & Una operación de \textit{diff} entre pricings que el motor no tiene \\
\hline
\multicolumn{4}{|l|}{\textbf{D. Relación entre límites de uso y \textit{rates} (RQ3) (6)}} \\
\hline
PI-D01 & En \texttt{CONFIGURATION} de \texttt{SAAS}, qué se agota antes, la cuota mensual o el \textit{rate} & P, demanda sostenible & \\
PI-D02 & A partir de qué volumen mensual el \textit{rate} de \texttt{SAAS} deja de ser el límite efectivo & P, demanda sostenible & \\
PI-D03 & Si \texttt{SAAS} cobra por correo y limita por petición, cuántos destinatarios por petición igualan los dos límites & P, cuotas que se agotan a la vez & \\
PI-D04 & Si el \textit{rate} de \texttt{SAAS} sube al subir de plan o es el mismo en todos & D & \\
PI-D05 & Si los límites de \texttt{SAAS} crecen con el uso (reputación) o son fijos por plan & D & \\
PI-D06 & Qué ocurre en \texttt{SAAS} al superar el \textit{rate}: rechazo, cola o degradación & X & El comportamiento al exceder no es una dimensión del ACM \\
\hline
\end{tabular}
\caption[Catálogo de preguntas de decisión, categorías C y D]{Catálogo de preguntas de decisión, categorías C y D. Misma
leyenda que la Tabla~\ref{tab:catalogo-pi-1}.}
\label{tab:catalogo-pi-2}
\end{table}

\begin{table}[htbp]
\centering
\scriptsize
\begin{tabular}{|l|>{\raggedright\arraybackslash}p{6.6cm}|>{\raggedright\arraybackslash}p{2.2cm}|>{\raggedright\arraybackslash}p{3.4cm}|}
\hline
\textbf{Id} & \textbf{Pregunta (plantilla abreviada)} & \textbf{Clase y operación} & \textbf{Qué falta (X) o cómo se lee (L)} \\
\hline
\multicolumn{4}{|l|}{\textbf{E. Necesidades funcionales (3)}} \\
\hline
PI-E01 & Si \texttt{CONFIGURATION} de \texttt{SAAS} incluye \texttt{FEATURE} & L & Por lectura de la matriz de \textit{features}. Con instrumento exigiría el iPricing del canal web, fuera del alcance fijado en el Capítulo~\ref{cap:catalogo} \\
PI-E02 & Cuál de \texttt{SAAS 1} y \texttt{SAAS 2} da \texttt{FEATURE} al menor precio para \texttt{USAGE LIMIT} correos al mes & L & Por lectura: \textit{feature} en la matriz, capacidad en la Datasheet, composición a mano. Con instrumento faltaría el vínculo entre canales \\
PI-E03 & Cuánto tiempo conserva \texttt{SAAS} los registros de envío en \texttt{CONFIGURATION} & X & Un límite de uso que no se consume (retención): v0.3 no lo modela \\
\hline
\multicolumn{4}{|l|}{\textbf{F. Evolución y vigencia (2)}} \\
\hline
PI-F01 & Si el pricing de \texttt{SAAS} ha cambiado desde \texttt{YEAR} en cuota o en \textit{rate} & D & \\
PI-F02 & Qué parte de la Datasheet de \texttt{SAAS} sigue vigente & X & Una Datasheet no dice qué parte sigue vigente (\mbox{F-09}, la caducidad de los datos): exige reverificar fuente a fuente \\
\hline
\multicolumn{4}{|l|}{\textbf{G. Operación y riesgo (3)}} \\
\hline
PI-G01 & Si el tráfico crece un \texttt{USAGE LIMIT}~\% al mes, cuándo habrá que cambiar de plan en \texttt{SAAS} & H, suscripción más barata iterada & \\
PI-G02 & Coste de un mes de \texttt{USAGE LIMIT} correos en \texttt{SAAS} sin cambiar de plan & P, importe del sobrecoste & \\
PI-G03 & Si se puede contratar más volumen en \texttt{SAAS} sin cambiar de plan, y a qué precio & D & \\
\hline
\end{tabular}
\caption[Catálogo de preguntas de decisión, categorías E a G]{Catálogo de preguntas de decisión, categorías E a G. Misma
leyenda que la Tabla~\ref{tab:catalogo-pi-1}.}
\label{tab:catalogo-pi-3}
\end{table}

Las cinco preguntas X no se descartan: cada una deja anotado en la última columna de las tablas lo que le falta al instrumental.
PI-E03, la retención de los registros de envío, pide un límite que no se consume con cada envío, y PI-D06, el comportamiento al superar el \textit{rate}, algo que ninguna dimensión de capacidad del ACM recoge.
Una pregunta X solo pasa a respondible cuando el esquema o el motor ganan lo que le falta, y por eso la cobertura del 78~\% se mide sobre el instrumental actual.

Solo una de las cinco tiene una forma esbozada en el Capítulo~\ref{cap:evolucion}: PI-F02, la vigencia de la Datasheet.
Se respondería con una fecha de verificación en cada afirmación del registro de procedencia, la forma que la Tabla~\ref{tab:sinforma} esboza para F-09, la caducidad de los datos.
P-09 (Apartado~\ref{subsec:evolucion:p09}) declara qué hace una cuota al superarse, la cola de reintento de Brevo, pero no cubre el exceso sobre el \textit{rate} que pide PI-D06.
La concurrencia de PI-A08, el \textit{diff} entre pricings de PI-C05 y la retención de PI-E03 no tienen extensión propuesta, y quedan como límite del esquema y del motor.
